\documentclass[aspectratio=169]{beamer}

\usetheme{Madrid}
\usecolortheme{default}
\setbeamertemplate{navigation symbols}{}
\setbeamertemplate{footline}{%
  \leavevmode%
  \hbox{%
    \begin{beamercolorbox}[wd=\paperwidth,ht=2.25ex,dp=1ex,right]{date in head/foot}%
      \usebeamerfont{date in head/foot}\insertframenumber{} / \inserttotalframenumber\hspace*{2ex}%
    \end{beamercolorbox}%
  }%
}

\usepackage[utf8]{inputenc}
\usepackage[T1]{fontenc}
\usepackage{amsmath,amssymb,mathtools}
\usepackage{booktabs}
\usepackage{graphicx}
\usepackage{array}

\title[Nesterov smoothing]{Nesterov's Smooth Minimization of Non-Smooth Functions}
\author{Micha\l{} Szewczak \and Sebastian Trojan}
\date{\today}

\newcommand{\maybeimage}[2][]{%
  \IfFileExists{#2}{\includegraphics[#1]{#2}}{%
    \IfFileExists{results/plots/#2}{\includegraphics[#1]{results/plots/#2}}{%
      \fbox{\begin{minipage}[c][0.35\textheight][c]{0.8\textwidth}\centering Missing image: \texttt{\detokenize{#2}}\end{minipage}}%
    }%
  }%
}

\begin{document}

% 1
\begin{frame}
  \titlepage
  \vspace{-1em}
  \begin{center}
    \small Based on Yurii Nesterov, \emph{Smooth minimization of non-smooth functions}, 2005.
  \end{center}
\end{frame}

% 2
\begin{frame}{Motivation}
  Many convex optimization problems are non-smooth because their objectives contain maxima, norms, or absolute values.

  \vspace{0.6em}
  Classical subgradient methods for non-smooth convex optimization have a worst-case rate of order
  \[
      O(1/\varepsilon^2).
  \]

  Nesterov's smoothing method uses additional max-structure of the objective. It replaces the non-smooth term by a smooth approximation and then applies an accelerated first-order method.
  The resulting theoretical complexity is improved to
  \[
      O(1/\varepsilon).
  \]

  \vspace{0.3em}
  The main trade-off is controlled by the smoothing parameter: smaller $\mu$ gives a more accurate approximation, but a less smooth problem.
\end{frame}

% 3
\begin{frame}{Structured problem class}
  The original problem is
  \[
      f^*=\min_{x\in Q_1}f(x),
  \]
  where
  \[
      f(x)=\hat f(x)+\max_{u\in Q_2}
      \{\langle Ax,u\rangle-\hat\phi(u)\}.
  \]

  \vspace{0.5em}
  \begin{itemize}
    \item $x$ is the primal variable.
    \item $u$ is the auxiliary, or dual, variable.
    \item $\hat f(x)$ is the smooth part of the objective.
    \item $\hat\phi(u)$ is the part of the max-representation depending only on $u$.
    \item In the implemented examples, $\hat f(x)=0$.
  \end{itemize}

  \vspace{0.3em}
  The non-smoothness comes from the maximization over $u$.
\end{frame}

% 4
\begin{frame}{Smoothing the max-term}
  Choose a strongly convex prox-function $d_2$ on $Q_2$, normalized so that
  \[
      \min_{u\in Q_2}d_2(u)=0,
      \qquad
      D_2=\max_{u\in Q_2}d_2(u).
  \]

  Define the smoothed max-term
  \[
      f_\mu(x)=\max_{u\in Q_2}
      \{\langle Ax,u\rangle-\hat\phi(u)-\mu d_2(u)\}.
  \]
  The full smoothed objective is
  \[
      \bar f_\mu(x)=\hat f(x)+f_\mu(x).
  \]

  The approximation error satisfies
  \[
      \bar f_\mu(x)\le f(x)\le \bar f_\mu(x)+\mu D_2.
  \]
  Therefore, choosing
  \[
      \mu=\frac{\varepsilon}{2D_2}
  \]
  makes the smoothing error at most $\varepsilon/2$.
\end{frame}

% 5
\begin{frame}{Gradient and smoothness trade-off}
  The smoothed dual response is
  \[
      u_\mu(x)=\arg\max_{u\in Q_2}
      \{\langle Ax,u\rangle-\hat\phi(u)-\mu d_2(u)\}.
  \]

  Then
  \[
      \nabla f_\mu(x)=A^*u_\mu(x),
      \qquad
      \nabla\bar f_\mu(x)=\nabla\hat f(x)+A^*u_\mu(x).
  \]

  The gradient of $\bar f_\mu$ is Lipschitz continuous with constant
  \[
      L_\mu=M+\frac{\|A\|_{1,2}^2}{\mu\sigma_2}.
  \]

  \vspace{0.3em}
  Here $M$ is the Lipschitz constant of $\nabla\hat f$, and $\|A\|_{1,2}$ is the operator norm associated with the chosen primal and dual geometries.

  \vspace{0.3em}
  Decreasing $\mu$ improves approximation accuracy, but increases $L_\mu$.
\end{frame}

% 6
\begin{frame}[t]{Accelerated minimization}
\scriptsize
For a fixed smoothing parameter $\mu$, the method minimizes
\[
    \min_{x\in Q_1}\bar f_\mu(x).
\]

\vspace{-0.5em}
\begin{block}{One accelerated iteration}
\[
\begin{array}{rcll}
 g_k &=& \nabla\bar f_\mu(x_k) & \text{gradient},\\[0.25em]
 y_k &=& \displaystyle\arg\min_{y\in Q_1}\left\{\langle g_k,y-x_k\rangle+\frac{L_\mu}{2}\|y-x_k\|^2\right\} & \text{local point},\\[0.45em]
 G_k &=& \displaystyle\sum_{i=0}^{k}\alpha_i g_i,\qquad \alpha_i=\frac{i+1}{2} & \text{accumulated gradient},\\[0.45em]
 z_k &=& \displaystyle\arg\min_{x\in Q_1}\left\{\frac{L_\mu}{\sigma_1}d_1(x)+\langle G_k,x\rangle\right\} & \text{aggregate point},\\[0.65em]
 x_{k+1} &=& \displaystyle\frac{2}{k+3}z_k+\frac{k+1}{k+3}y_k & \text{next search point}.
\end{array}
\]
\end{block}

\vspace{-0.2em}
$y_k$ uses the current gradient, while $z_k$ uses all gradients accumulated in $G_k$.
\end{frame}

% 7
\begin{frame}{Primal-dual certificate}
\small
  The algorithm builds an averaged dual point
  \[
      \hat u=\frac{\sum_{k=0}^{N}\alpha_k u_k}{\sum_{k=0}^{N}\alpha_k},
      \qquad u_k=u_\mu(x_k),
      \qquad \hat x=y_N.
  \]

  The dual lower-bound function is
  \[
      \phi(u)=\min_{x\in Q_1}
      \{\hat f(x)+\langle Ax,u\rangle\}-\hat\phi(u).
  \]


  \vspace{0.3em}
  By weak duality,
  \[
      \phi(\hat u)\le f^*\le f(\hat x).
  \]
  Therefore the primal-dual gap
  \[
      f(\hat x)-\phi(\hat u)
  \]
  certifies the solution quality.
\end{frame}

% 8
\begin{frame}{Theoretical iteration bound}
\small
  Nesterov's estimate gives
  \[
      f(\hat x)-\phi(\hat u)
      \le
      \frac{4\|A\|_{1,2}}{N+1}
      \sqrt{\frac{D_1D_2}{\sigma_1\sigma_2}}
      +
      \frac{4MD_1}{\sigma_1(N+1)^2}.
  \]

  In the implemented examples, $\hat f=0$, so $M=0$. Hence
  \[
      f(\hat x)-\phi(\hat u)
      \le
      \frac{4\|A\|_{1,2}}{N+1}
      \sqrt{\frac{D_1D_2}{\sigma_1\sigma_2}}.
  \]

  To guarantee $f(\hat x)-\phi(\hat u)\le\varepsilon$, it is sufficient to take
  \[
      N_{\mathrm{pred}}+1=
      \left\lceil
      \frac{4\|A\|_{1,2}}{\varepsilon}
      \sqrt{\frac{D_1D_2}{\sigma_1\sigma_2}}
      \right\rceil.
  \]

  Thus the predicted complexity is
  \[
      N=O(1/\varepsilon).
  \]
\end{frame}

% 9
\begin{frame}{Problem classes}
\small
  \[
  \textbf{Matrix game:}\qquad
  \min_{x\in\Delta_n}\max_{u\in\Delta_m}u^\top Ax
  \]

  \[
  \textbf{Continuous location:}\qquad
  \min_{\|x\|_2\le r}\sum_{j=1}^{p}w_j\|x-c_j\|_2
  \]

  \[
  \textbf{Maximum of absolute values:}\qquad
  \min_{\|x\|_2\le r}\max_{1\le j\le m}|a_j^\top x-b_j|
  \]

  \[
  \textbf{Sum of absolute values:}\qquad
  \min_{\|x\|_2\le r}\sum_{j=1}^{m}|a_j^\top x-b_j|
  \]
\end{frame}

% 10
\begin{frame}{Experimental plan}
  \begin{enumerate}
    \item Reproduce Nesterov's matrix game experiment.
    \item Apply the same scheme to the other three examples.
    \item Test the empirical dependence on $\varepsilon$.
    \item Compare standard smoothing with a changing-$\mu$ strategy.
  \end{enumerate}

  \vspace{0.6em}
  In every experiment, convergence is measured by the certified condition
  \[
      f(\hat x)-\phi(\hat u)\le\varepsilon.
  \]
\end{frame}

% 11
\begin{frame}{Experiment 1: matrix game reproduction}
  Reproduced problem:
  \[
      \min_{x\in\Delta_n}\max_{u\in\Delta_m}u^\top Ax,
      \qquad A_{ij}\sim U[-1,1].
  \]

  \vspace{0.5em}
  Tested grid:
  \begin{itemize}
    \item $\varepsilon\in\{10^{-2},10^{-3},10^{-4}\}$,
    \item $m\in\{100,300,1000\}$,
    \item $n\in\{100,300,1000,3000,10000\}$, with the largest $n$ omitted for $\varepsilon=10^{-4}$.
  \end{itemize}

  \vspace{0.5em}
  Each table cell reports iterations / percentage of the predicted iteration bound.
\end{frame}

% 12
\begin{frame}{Experiment 1: matrix game results}
\tiny
\centering
$\varepsilon=10^{-2}$\vspace{0.1em}

\resizebox{0.70\textwidth}{!}{%
\begin{tabular}{lccccc}
\toprule
$m\backslash n$ & 100 & 300 & 1\,000 & 3\,000 & 10\,000\\
\midrule
100 & 765/42\% & 934/46\% & 1\,108/49\% & 1\,187/49\% & 1\,281/49\%\\
300 & 690/34\% & 910/40\% & 1\,126/45\% & 1\,268/47\% & 1\,475/51\%\\
1\,000 & 761/34\% & 873/35\% & 1\,086/39\% & 1\,284/43\% & 1\,457/46\%\\
\bottomrule
\end{tabular}}

\vspace{0.35em}
$\varepsilon=10^{-3}$\vspace{0.1em}

\resizebox{0.70\textwidth}{!}{%
\begin{tabular}{lccccc}
\toprule
$m\backslash n$ & 100 & 300 & 1\,000 & 3\,000 & 10\,000\\
\midrule
100 & 7\,123/39\% & 8\,426/41\% & 8\,928/40\% & 10\,191/42\% & 10\,080/39\%\\
300 & 7\,337/36\% & 8\,639/38\% & 10\,154/40\% & 11\,302/42\% & 12\,022/41\%\\
1\,000 & 7\,526/33\% & 8\,828/35\% & 10\,157/37\% & 11\,595/39\% & 14\,167/44\%\\
\bottomrule
\end{tabular}}

\vspace{0.35em}
$\varepsilon=10^{-4}$\vspace{0.1em}

\resizebox{0.60\textwidth}{!}{%
\begin{tabular}{lcccc}
\toprule
$m\backslash n$ & 100 & 300 & 1\,000 & 3\,000\\
\midrule
100 & 68\,505/37\% & 64\,512/31\% & 77\,136/34\% & 74\,024/30\%\\
300 & 60\,932/30\% & 78\,680/34\% & 83\,844/33\% & 92\,862/34\%\\
1\,000 & 70\,778/31\% & 77\,644/31\% & 93\,346/34\% & 102\,925/35\%\\
\bottomrule
\end{tabular}}

\end{frame}

% 13
\begin{frame}{Experiment 2: other problem classes}
\small
  The matrix game experiment was extended to:
  \begin{itemize}
    \item Continuous location:
    \(\min_{\|x\|_2\le r}\sum_{j=1}^p w_j\|x-c_j\|_2,\quad c_j\sim \mathrm{Unif}(B_2(0,0.8r)),\ w_j\sim U[0.5,1.5],\ r=1\).\par
    Grid: \(p\in\{10,50\}\), \(d\in\{2,5\}\) where \(c_j,x\in\mathbb{R}^d\).
    \item Maximum of absolute values:
    \(\min_{\|x\|_2\le r}\max_{1\le j\le m}|a_j^\top x-b_j|,\quad A_{ji}, b_j\sim U[-1,1],\ r=1\).\par
    Grid: \(m\in\{100,300\}\), \(n\in\{50,200\}\) where \(a_j,x\in\mathbb{R}^n\).
    \item Sum of absolute values:
    \(\min_{\|x\|_2\le r}\sum_{j=1}^m|a_j^\top x-b_j|,\quad A_{ji}, b_j\sim U[-1,1],\ r=1\).\par
    Grid: \(m\in\{100,300\}\), \(n\in\{50,200\}\) where \(a_j,x\in\mathbb{R}^n\).
  \end{itemize}

  \vspace{0.5em}
  Common accuracy levels:
  \[
      \varepsilon\in\{10^{-2},10^{-3},10^{-4}\}.
  \]

\end{frame}

% 14
\begin{frame}{Experiment 2: summary tables}
\tiny
\centering
Continuous location\vspace{0.1em}

\resizebox{0.54\textwidth}{!}{%
\begin{tabular}{lrrrr}
\toprule
$\varepsilon$ & Avg. iter. & Max iter. & Avg. \% & Time\\
\midrule
$10^{-4}$ & 171\,599 & 320\,867 & 30.64\% & 31.8s\\
$10^{-3}$ & 20\,525 & 31\,653 & 43.74\% & 3.86s\\
$10^{-2}$ & 1\,988 & 3\,606 & 36.70\% & 0.35s\\
\bottomrule
\end{tabular}}

\vspace{0.15em}
Maximum of absolute values\vspace{0.1em}

\resizebox{0.54\textwidth}{!}{%
\begin{tabular}{lrrrr}
\toprule
$\varepsilon$ & Avg. iter. & Max iter. & Avg. \% & Time\\
\midrule
$10^{-4}$ & 171\,486 & 249\,861 & 39.81\% & 10.6s\\
$10^{-3}$ & 20\,032 & 33\,671 & 44.76\% & 1.26s\\
$10^{-2}$ & 2\,593 & 3\,788 & 56.06\% & 0.16s\\
\bottomrule
\end{tabular}}

\vspace{0.15em}
Sum of absolute values\vspace{0.1em}

\resizebox{0.54\textwidth}{!}{%
\begin{tabular}{lrrrr}
\toprule
$\varepsilon$ & Avg. iter. & Max iter. & Avg. \% & Time\\
\midrule
$10^{-4}$ & 1\,198\,404 & 2\,711\,368 & 29.45\% & 54.9s\\
$10^{-3}$ & 135\,894 & 284\,262 & 32.52\% & 6.43s\\
$10^{-2}$ & 12\,088 & 27\,750 & 30.43\% & 0.51s\\
\bottomrule
\end{tabular}}

\end{frame}

% 15
\begin{frame}{Experiment 3: testing the $O(1/\varepsilon)$ behavior}
  To isolate the dependence on $\varepsilon$, five seeds were used; for each seed, one instance was generated and then kept fixed across all values of $\varepsilon$:
  \[
      \varepsilon\in\{10^{-1},5\cdot10^{-2},10^{-2},5\cdot10^{-3},10^{-3},5\cdot10^{-4},10^{-4}\}.
  \]

  \vspace{0.4em}
  The theory predicts
  \[
      N(\varepsilon)\approx C\frac{1}{\varepsilon}.
  \]

  \vspace{0.4em}
  Therefore:
  \begin{itemize}
    \item iterations versus $1/\varepsilon$ should be approximately linear,
    \item the log-log fitted exponent $\hat p$ should be close to $1$.
  \end{itemize}
\end{frame}

% 16
\begin{frame}{Experiment 3: iterations versus $1/\varepsilon$}
  \centering
  \maybeimage[height=0.70\textheight]{epsilon_scaling_iterations_vs_inverse_epsilon.png}

  \vspace{0.2em}
  \small The curves are close to linear; the difference between problems is mainly the constant factor.
\end{frame}

% 17
\begin{frame}{Experiment 3: log-log scaling}
\begin{columns}[T,onlytextwidth]
  \begin{column}{0.58\textwidth}
    \centering
    \maybeimage[width=\linewidth]{epsilon_scaling_loglog.png}
  \end{column}
  \begin{column}{0.39\textwidth}
    \small
    Empirical power-law model:
    \[
        N(\varepsilon)\approx C\varepsilon^{-p}.
    \]
    Taking logarithms gives
    \[
        \log N=\log C+p\log(1/\varepsilon).
    \]
    Therefore $\hat p$ is the slope of the fitted line on the log-log plot. The theory predicts
    \[
        p=1.
    \]

    \vspace{0.2em}
    \scriptsize
    \begin{tabular}{lcc}
    \toprule
    Problem & $\hat p$ & SE\\
    \midrule
    Matrix game & 0.9897 & 0.0096\\
    Continuous location & 0.9986 & 0.0054\\
    Max. abs. values & 0.9208 & 0.0042\\
    Sum abs. values & 0.9649 & 0.0096\\
    \bottomrule
    \end{tabular}
  \end{column}
\end{columns}
\end{frame}

% 18
\begin{frame}{Experiment 4: changing the smoothing parameter}
\footnotesize
\textbf{1. Motivation}
\begin{itemize}
  \item Standard smoothing uses $\mu_\varepsilon=\varepsilon/(2D_2)$ from the beginning.
  \item For small $\varepsilon$, this value is small, so the smoothed problem is sharper and harder.
\end{itemize}

\vspace{0.2em}
\textbf{2. Method}
\begin{itemize}
  \item Start with a larger smoothing parameter and decrease it geometrically:
  \[
      \mu_0=c\mu_\varepsilon,
      \qquad
      \mu_{t+1}=\rho\mu_t,
      \qquad
      c>1,\ \rho\in(0,1).
  \]
  \item Move to the next stage when
  \[
      f(\hat x)-\phi(\hat u)\le c_{\mathrm{stage}}\mu_tD_2,
      \qquad
      c_{\mathrm{stage}}\ge 1.
  \]
\end{itemize}

\vspace{0.2em}
\textbf{3. Our parameters}
\[
    c=8,
    \qquad
    \rho=0.5,
    \qquad
    c_{\mathrm{stage}}=1,
    \qquad
    8\mu_\varepsilon\to4\mu_\varepsilon\to2\mu_\varepsilon\to\mu_\varepsilon.
\]
\end{frame}

% 19
\begin{frame}{Experiment 4: gap history}
  \centering
  \maybeimage[height=0.70\textheight]{mu_gap_history_matrix_game_eps_1e-04.png}

  \vspace{0.2em}
  \small After reducing $\mu$, the smoothed objective changes and the primal-dual gap can jump upward. In most tested cases, continuation required more iterations.
\end{frame}

% 20
\begin{frame}{Conclusions}
  \begin{itemize}
    \item The implementation reached the required primal-dual gap on all reported instances.
    \item Matrix game results reproduce the qualitative behavior from Nesterov's paper.
    \item The empirical scaling is consistent with the theoretical $O(1/\varepsilon)$ complexity.
    \item Worst-case iteration bounds were conservative on random instances.
    \item Practical difficulty depends strongly on the problem structure.
    \item The tested changing-$\mu$ strategy was usually less efficient than the standard smoothing method.
  \end{itemize}
\end{frame}

% 21
\begin{frame}{Reference}
  \small
  \begin{thebibliography}{1}
  \bibitem{nesterov2005}
  Yurii Nesterov.
  \newblock \emph{Smooth minimization of non-smooth functions}.
  \newblock Mathematical Programming, 103(1):127--152, 2005.
  \end{thebibliography}
\end{frame}

\end{document}
